\chapter{Simulation Experiments}
\label{ch:simulations}

This chapter validates the theoretical results of Chapter~\ref{ch:convergence} through numerical experiments on two-player matrix games. Section~\ref{sec:ewpg-experiments} tests the AM-HM variance improvement of evidence-weighted PG (Theorem~\ref{thm:am-hm-ch7}). Section~\ref{sec:lola-experiments} tests the basin-of-attraction enlargement of LOLA-PG (Theorem~\ref{thm:basin-ch7}).

All experiments use the REINFORCE gradient estimator (Algorithm~1 of~\cite{giannou2022}), corresponding to Model~3 in their framework. Step sizes follow $\gamma_n = \gamma/(n+m)^p$ with $p = 0.75$, $\gamma = 0.5$, $m = 10$. Policies are projected onto the probability simplex after each update. Source code is available at \texttt{simulations/evidence\_weighted\_pg.py} and \texttt{simulations/lola\_basin.py}.


%% ========================================================================
\section{Evidence-Weighted PG Experiments}
\label{sec:ewpg-experiments}
%% ========================================================================

\subsection{Setup}

We compare standard PG and EWPG on two games with known interior Nash equilibria:

\begin{center}
\renewcommand{\arraystretch}{1.2}
\begin{tabular}{lcccc}
\hline
\textbf{Game} & \textbf{Players} & \textbf{Actions} & \textbf{Nash} & \textbf{Type} \\
\hline
Matching Pennies & 2 & $2 \times 2$ & $(0.5, 0.5)$ & Zero-sum \\
Rock-Paper-Scissors & 2 & $3 \times 3$ & $(1/3, 1/3, 1/3)$ & Zero-sum \\
\hline
\end{tabular}
\end{center}

To simulate heterogeneous evidence quality, we set $V = (V_1, V_2) = (10, 1)$, giving agent~1 high-quality data and agent~2 low-quality data. Each agent's gradient estimate is corrupted by additive Gaussian noise with standard deviation $\sigma_i = C/\sqrt{V_i}$ where $C = 2$. Each experiment averages over 50 independent runs from a shared initial condition near the Nash equilibrium.

For EWPG, the evidence weights are $w_i = V_i / \max(V)$, so $w_1 = 1$ and $w_2 = 0.1$.

\subsection{Experiment 1: Convergence comparison}

Figure~\ref{fig:convergence-ewpg} plots the squared distance to Nash, $\|\pi_n - \pi^*\|^2$, on a logarithmic scale against the episode number. Both standard PG and EWPG converge, but EWPG exhibits lower variance in the distance trajectory, as predicted by Theorem~\ref{thm:am-hm-ch7}.

\begin{figure}[H]
\centering
\includegraphics[width=0.48\textwidth]{figures/ewpg/convergence_matching_pennies.png}
\hfill
\includegraphics[width=0.48\textwidth]{figures/ewpg/convergence_rock-paper-scissors.png}
\caption{Convergence of standard PG (red) vs.\ EWPG (green) on Matching Pennies (left) and Rock-Paper-Scissors (right). Evidence weights $V = (10, 1)$. Shaded regions show $\pm 2$ standard errors over 50 runs. Both converge at comparable rates, but EWPG maintains lower variance throughout, consistent with the HM/AM improvement.}
\label{fig:convergence-ewpg}
\end{figure}

\subsection{Experiment 2: Variance ratio vs.\ evidence heterogeneity}

To test the quantitative prediction of Theorem~\ref{thm:am-hm-ch7}, we sweep the evidence ratio $r = V_1/V_2$ over $\{1, 2, 3, 5, 8, 10, 15, 20\}$ and measure the empirical variance ratio $\sigma_w^2 / \sigma_{\mathrm{std}}^2$ from the last 500 episodes of each run.

The theoretical prediction is:
\begin{equation}
\label{eq:theoretical-ratio}
    \frac{\operatorname{HM}(V)}{\operatorname{AM}(V)} = \frac{4r}{(1+r)^2},
\end{equation}
which equals 1 at $r = 1$ (homogeneous agents) and decreases monotonically as $r$ increases.

\begin{figure}[H]
\centering
\includegraphics[width=0.48\textwidth]{figures/ewpg/variance_ratio_matching_pennies.png}
\hfill
\includegraphics[width=0.48\textwidth]{figures/ewpg/variance_ratio_rock-paper-scissors.png}
\caption{Empirical variance ratio (blue squares) vs.\ theoretical HM/AM prediction (black circles) as a function of evidence ratio $r = V_1/V_2$. Left: Matching Pennies. Right: Rock-Paper-Scissors. The empirical ratios track the theoretical curve closely, confirming the AM-HM variance improvement. At $r = 10$, the variance reduction is approximately 67\%.}
\label{fig:variance-ratio}
\end{figure}

Figure~\ref{fig:variance-ratio} shows close agreement between theory and experiment. At $r = 10$, the theoretical ratio is $4 \cdot 10 / 121 \approx 0.33$, and the empirical ratio is comparably small. The slight deviation at extreme ratios ($r = 15, 20$) is attributable to the increased difficulty of the learning problem when one agent is heavily down-weighted.


%% ========================================================================
\section{LOLA-PG Basin of Attraction Experiments}
\label{sec:lola-experiments}
%% ========================================================================

\subsection{Setup}

We compare the basin of attraction of standard PG and LOLA-PG on two $2 \times 2$ games:

\begin{center}
\renewcommand{\arraystretch}{1.2}
\begin{tabular}{lccc}
\hline
\textbf{Game} & \textbf{Nash $(p_1^*, p_2^*)$} & \textbf{Type} & \textbf{Spectral reinforcement?} \\
\hline
Matching Pennies & $(0.5, 0.5)$ & Zero-sum & Yes \\
Battle of the Sexes & $(0.6, 0.4)$ & Coordination & Partial \\
\hline
\end{tabular}
\end{center}

For $2 \times 2$ games, policies reduce to scalars $(p_1, p_2) \in [0,1]^2$, where $p_i$ is the probability of playing action~0. This allows us to visualise the entire policy space as a unit square.

We sweep initial conditions over a $25 \times 25$ grid in $[0.05, 0.95]^2$. For each initial condition, we run 8 independent trials of 1500 episodes with Gaussian noise ($\sigma = 0.3$) and declare convergence if $\|\pi_{\mathrm{final}} - \pi^*\| < 0.08$. LOLA-PG uses $\lambda = 0.5$ and $\eta = 0.1$.

\subsection{Experiment 3: Basin mapping}

Figure~\ref{fig:basin} shows heat maps of convergence probability across the policy space. Green indicates reliable convergence; red indicates divergence.

\begin{figure}[H]
\centering
\includegraphics[width=0.48\textwidth]{figures/lola/basin_matching_pennies.png}
\includegraphics[width=0.48\textwidth]{figures/lola/basin_battle_of_the_sexes.png}
\caption{Basin of attraction for standard PG (left panel), LOLA-PG (centre panel), and their difference (right panel) on Matching Pennies (top) and Battle of the Sexes (bottom). The white star marks the Nash equilibrium. LOLA-PG shows a strictly larger basin of convergence in both games, with the enlargement particularly pronounced in Matching Pennies (a zero-sum game where the spectral reinforcement condition holds).}
\label{fig:basin}
\end{figure}

For Matching Pennies, LOLA-PG approximately doubles the basin area compared to standard PG. This is consistent with Theorem~\ref{thm:basin-ch7}: in zero-sum games, the opponent-shaping Hessian converts rotational dynamics (which cause standard PG to orbit the Nash equilibrium) into contractive dynamics. The improvement is less dramatic in Battle of the Sexes, where the game is not purely zero-sum and the spectral reinforcement condition holds only partially.

\subsection{Experiment 4: Trajectory visualisation}

Figure~\ref{fig:trajectories} overlays sample trajectories from random initial conditions, coloured by convergence outcome.

\begin{figure}[H]
\centering
\includegraphics[width=0.48\textwidth]{figures/lola/trajectories_matching_pennies.png}
\hfill
\includegraphics[width=0.48\textwidth]{figures/lola/trajectories_battle_of_the_sexes.png}
\caption{Sample trajectories for standard PG (left) and LOLA-PG (right) on Matching Pennies (left pair) and Battle of the Sexes (right pair). Green trajectories converge to Nash; red trajectories diverge. Under standard PG, Matching Pennies trajectories tend to spiral around the Nash equilibrium without converging. LOLA-PG converts these spirals into convergent paths.}
\label{fig:trajectories}
\end{figure}

The trajectory plots reveal the mechanism behind the basin enlargement. In Matching Pennies under standard PG, the gradient field has a strong rotational component at the Nash equilibrium, causing policies to orbit rather than converge. The opponent-shaping term in LOLA-PG counteracts this rotation, converting the dynamics from circular to spiralling inward. This is precisely the effect captured by the spectral reinforcement condition: the symmetric part of the opponent-shaping Hessian adds negative-definite contraction that dominates the antisymmetric rotation.


%% ========================================================================
\section{Discussion}
\label{sec:sim-discussion}
%% ========================================================================

The experiments confirm both theoretical predictions:

\begin{enumerate}
    \item \textbf{EWPG:} The empirical variance ratio closely tracks the theoretical $\operatorname{HM}(V)/\operatorname{AM}(V)$ curve across a range of evidence heterogeneity levels. The improvement is most significant when agents have highly unequal data quality ($r \gg 1$).

    \item \textbf{LOLA-PG:} The basin of attraction is strictly larger under LOLA-PG than standard PG, with the enlargement most pronounced in zero-sum games where the spectral reinforcement condition is satisfied.
\end{enumerate}

\paragraph{Limitations of the experiments.}
\begin{itemize}
    \item All games tested are single-state ($|\mathcal{S}| = 1$), i.e., normal-form games rather than stochastic games. The theoretical results apply to the full stochastic game setting, but the experiments do not test this generality.
    \item The noise model (additive Gaussian) is a simplification of the REINFORCE estimator's actual variance structure, which depends on the policy through the minimum action probability $\kappa_i = \min_{s,a} \pi_i(a|s)$.
    \item The LOLA opponent-shaping term requires knowledge of the game's payoff structure to compute the cross-derivatives $\partial^2 V_i / \partial \pi_i \partial \pi_j$. In practice, these must be estimated, introducing additional error not captured in our analysis.
\end{itemize}

\paragraph{Future experiments.} Testing on multi-state stochastic games (e.g., grid worlds with strategic interaction, the algorithmic pricing environment of Calvano et al.~(2020)) and with the full REINFORCE estimator rather than additive noise would provide stronger empirical support.
